\section{Frontend and Abstract Syntax Tree Construction}\label{sec:implementation-frontend}

The frontend validates syntactic structure and builds the logical model of the composition.
It is implemented as a classic scanner--parser pipeline within the \texttt{motivo\_parsing} CMake target, with AST node
definitions shared from \texttt{common/ast/} headers consumed by later stages.

\subsection{Entry Points and Result Type}\label{detail-subsec:frontend-entry-points}

\texttt{parsing/parse.cpp} exposes two public functions declared in \texttt{parsing/parse.hpp}:

\noindent \textbf{\texttt{parse\_stream(FILE*, source\_name, diagnostics)}.} Used by \texttt{motivo::compile}; wraps
the input file in a \texttt{ScannerInputGuard} that resets Flex state, attaches the shared
\texttt{DiagnosticsEngine}, and calls \texttt{yyrestart}.

\noindent \textbf{\texttt{parse\_source(string, source\_name, diagnostics)}.} Used in unit tests; scans an in-memory
string via \texttt{yy\_scan\_string}.

Both paths construct an empty \texttt{ast::Program}, invoke \texttt{detail::Parser(loc, diagnostics,
program).parse()}, and return a \texttt{ParseResult} owning the completed tree.
If Bison returns \texttt{EXIT\_FAILURE} without prior diagnostics, \texttt{parse.cpp} emits a generic parse error.

\subsection{Lexical Analysis: Identifying Musical
Literals}\label{detail-subsec:lexical-analysis:-identifying-musical-literals}

The scanner is defined in \texttt{parsing/detail/scanner.l} and generated by Flex at build time.
It transforms the character stream into Bison token values (\texttt{Parser::symbol\_type}) while maintaining a
\texttt{source::Location} per token via \texttt{YY\_USER\_ACTION} column/line updates.

\subsubsection{Pitch-Class and Note Recognition}

Unlike conventional languages that treat musical data as strings, Motivo recognizes pitch literals such as
\texttt{F\#4} as primitive constants.
The scanner uses dedicated character-class macros:

\begin{lstlisting}[caption={Flex regular expression macros for note literal recognition (\texttt{scanner.l}).},label={detail-lst:scanner-note-macros}]
pitch        [A-G]
accidental   [#b]
octave       [0-8]
note_lit     {pitch}{accidental}?{octave}
\end{lstlisting}

When \texttt{note\_lit} matches, the action calls \texttt{parse\_note\_literal} from
\texttt{parsing/detail/scanner\_helpers.cpp}, which constructs a \texttt{music::Note} (pitch enum, accidental enum,
octave integer) and returns a \texttt{TOK\_NOTE\_LIT} token carrying that value.
Drum literals (\texttt{kick}, \texttt{snare}, etc.) and instrument keywords (\texttt{piano}, \texttt{drums}) are
recognized in the same file and mapped to \texttt{music::DrumNote} and \texttt{music::Instrument} enums.

\subsubsection{Keyword and Symbol Disambiguation}

Reserved keywords include structural tokens (\texttt{track}, \texttt{pattern}, \texttt{voice}, \texttt{play}) and
type keywords (\texttt{int}, \texttt{double}, \texttt{bool}, \texttt{note}, \texttt{seq}, \texttt{chord}) used in
\texttt{VarDeclStatement} and \texttt{TypedParameter} positions.
Identifiers match \texttt{\{alpha\}\{alnum\}*}; because keywords are listed before the identifier rule, longest-match
disambiguation ensures names such as \texttt{loop\_count} are not tokenized as the \texttt{loop} keyword.

Lexical errors call \texttt{report\_lexical}, which forwards to \texttt{DiagnosticsEngine::report} with
\texttt{DiagnosticStage::Parsing} when a diagnostics pointer is set by \texttt{scanner\_set\_diagnostics}.

\subsection{Syntactic Analysis: Building the Hierarchy}\label{detail-subsec:syntactic-analysis:-building-the-hierarchy}

The parser grammar lives in \texttt{parsing/detail/parser.y} and is processed by Bison~3.8 with C++ variant-based
semantic values (\texttt{\%define api.value.type variant}).
The generated \texttt{motivo::parsing::detail::Parser} class receives three parse parameters: a
\texttt{source::Location}, a \texttt{DiagnosticsEngine\&}, and the \texttt{ast::Program\&} under construction.

Grammar actions use helper macros \texttt{expression(...)} and \texttt{statement(...)} that allocate
\texttt{std::unique\_ptr\textless ast::Expression\textgreater{}} and \texttt{std::unique\_ptr\textless
ast::Statement\textgreater{}} nodes.
The start symbol \texttt{program} assembles \texttt{program.header}, \texttt{program.globals}, and
\texttt{program.tracks} from nested \texttt{track\_def}, \texttt{pattern\_def}, and \texttt{voice\_def} productions.
Detailed productions mirror the EBNF documented in \texttt{compiler/docs/grammar.bnf}.

\subsubsection{AST Node Architecture and Sum-Types}

The AST lives in \texttt{motivo::ast} headers under \texttt{common/ast/}.
Polymorphic nodes use \texttt{std::variant} for extensibility without virtual dispatch.

\paragraph{Declaration and Definition Nodes}
The top-level structure of the AST is represented by definitions that manage scoping and vertical layering:

\noindent \textbf{\texttt{TrackDefinition}:} Optional name, optional \texttt{music::Instrument}, and a body of
\texttt{TrackItem} variant nodes (statement, local pattern, or voice).

\noindent \textbf{\texttt{PatternDefinition}:} Name, \texttt{std::vector\{TypedParameter\}} (each holding
\texttt{types::Type} and identifier), and a statement block.

\noindent \textbf{\texttt{VoiceDefinition}:} Optional \texttt{from} expression and a body of \texttt{VoiceItem} nodes.

\paragraph{Statement and Expression Polymorphism}
Leaf nodes divide into statements (actions) and expressions (values):

\noindent \textbf{\texttt{StatementKind}:} \texttt{VarDeclStatement}, \texttt{AssignStatement}, \texttt{ForStatement},
\texttt{IfStatement}, \texttt{LoopStatement}, and \texttt{PlayStatement}.
The \texttt{VarDeclStatement} stores the declared \texttt{types::Type}, identifier, and initializer; it replaces
version~1.0's untyped \texttt{LetStatement}.

\noindent \textbf{\texttt{ExpressionKind}:} Literals (\texttt{Int}, \texttt{Float}, \texttt{Bool}, \texttt{Note},
\texttt{Rest}, \texttt{DrumNote}), \texttt{IdentifierExpression}, unary/binary/ternary operators,
\texttt{SequenceExpression}, \texttt{ChordExpression}, \texttt{PatternCallExpression}, and parenthesised wrappers.

\subsubsection{Recursive Node Ownership}

Recursive subtrees use \texttt{std::unique\_ptr} (e.g., \texttt{StatementPtr}, \texttt{ExpressionPtr}) so the AST
remains a tree in memory.
Every node carries a \texttt{source::Location} copied from the Bison \texttt{loc} parameter at reduction time, which
later drives diagnostic messages in semantic analysis and lowering failures.

The combination of variant dispatch and pointer-based recursion provides the foundation for the single-pass semantic
traversal and the interpreter-style lowering pass that follow.
